\documentclass[format=acmsmall, review=false]{acmart}
\usepackage{acm-ec-26}

% acmart already loads amsmath, amsthm (theorem, proposition, lemma, corollary, definition, example),
% graphicx, booktabs, hyperref, microtype and natbib.
\usepackage{mathtools}
\usepackage{subcaption}
\usepackage[capitalize,noabbrev]{cleveref}
\usepackage[textsize=footnotesize]{todonotes}
\presetkeys{todonotes}{inline}{}  % margins are narrow in acmsmall: all to-do notes inline

\setcitestyle{authoryear}

% numbers quoted in the text: generated by scripts/paper_tables.py from results/tables/v2
% generated by scripts/paper_tables.py -- do not edit
\newcommand{\bootRepsShapley}{50}
\newcommand{\effReal}{839}
\newcommand{\extFitHold}{6.3}
\newcommand{\extFitIn}{7.0}
\newcommand{\extHubs}{118}
\newcommand{\extQ}{7.3}
\newcommand{\feHoldQual}{36.2}
\newcommand{\feInQual}{46.0}
\newcommand{\feShare}{42.1}
\newcommand{\feShareHold}{32.3}
\newcommand{\holdGain}{6.62}
\newcommand{\holdGainFE}{6.632}
\newcommand{\holdGainNoFE}{6.625}
\newcommand{\holdTopOne}{14.2}
\newcommand{\holdTopTen}{39.8}
\newcommand{\idDeg}{-0.08}
\newcommand{\idDegSe}{0.10}
\newcommand{\idHigher}{48.6}
\newcommand{\idN}{638}
\newcommand{\idRec}{0.84}
\newcommand{\idRecSe}{0.08}
\newcommand{\jointDonor}{2.4}
\newcommand{\jointOwn}{4.5}
\newcommand{\logdegPps}{0.344}
\newcommand{\logdegUnif}{0.435}
\newcommand{\lrInertia}{44,813}
\newcommand{\lrMatch}{15,769}
\newcommand{\lrPop}{1,385,594}
\newcommand{\lrQual}{25}
\newcommand{\marketMedian}{344,678}
\newcommand{\mrsOrg}{6.6}
\newcommand{\mrsRec}{8.5}
\newcommand{\mrsTask}{246}
\newcommand{\mrsTaskHi}{407}
\newcommand{\mrsTaskLo}{149}
\newcommand{\nAdapter}{243,643}
\newcommand{\nAuthorsAll}{391,746}
\newcommand{\nChoices}{422,714}
\newcommand{\nDecisions}{199,464}
\newcommand{\nEdges}{574,531}
\newcommand{\nFinetune}{188,819}
\newcommand{\nFirst}{133,390}
\newcommand{\nFirstAuthors}{69,763}
\newcommand{\nFirstParents}{15,450}
\newcommand{\nModels}{1,860,413}
\newcommand{\nTrain}{117,509}
\newcommand{\placeboAlone}{3.1}
\newcommand{\recFullInertia}{3.2}
\newcommand{\recFullMatch}{5.2}
\newcommand{\recFullPop}{91.6}
\newcommand{\recQfFit}{94.0}
\newcommand{\recSfPop}{96.9}
\newcommand{\shHoldInertia}{3.6}
\newcommand{\shHoldMatch}{4.6}
\newcommand{\shHoldPopularity}{91.8}
\newcommand{\shHoldQuality}{0.0}
\newcommand{\shInertia}{3.2}
\newcommand{\shInertiaHi}{3.6}
\newcommand{\shInertiaLo}{2.7}
\newcommand{\shMatch}{5.2}
\newcommand{\shMatchHi}{5.4}
\newcommand{\shMatchLo}{5.0}
\newcommand{\shPopularity}{91.6}
\newcommand{\shPopularityHi}{92.0}
\newcommand{\shPopularityLo}{91.1}
\newcommand{\shQuality}{0.0}
\newcommand{\shQualityHi}{0.0}
\newcommand{\shQualityLo}{0.0}
\newcommand{\shareParentZero}{11}
\newcommand{\shareText}{75}
\newcommand{\sigMatch}{8.3}
\newcommand{\sigMatchRw}{6.4}
\newcommand{\sigPop}{87.3}
\newcommand{\sigShare}{66}
\newcommand{\tEnd}{July 2025}
\newcommand{\tHold}{May 5, 2025}
\newcommand{\tStart}{March 2022}
\newcommand{\tieMean}{7.9}
\newcommand{\topOneReal}{6.8}
\newcommand{\topOneUnif}{93.1}
\newcommand{\topTenReal}{21.9}
\newcommand{\valIn}{19}
\newcommand{\valN}{25}
\newcommand{\valSeRatio}{0.65}


\newcommand{\E}{\mathbb{E}}
\newcommand{\1}{\mathbf{1}}
\newcommand{\D}{\mathcal{D}}
\newcommand{\M}{\mathcal{M}}

\title[Popularity or Fit?]{Popularity or Fit? A Discrete-Choice Decomposition of Base-Model Selection on Hugging Face}

% Anonymized submission: no names, no acknowledgements.
\author{Anonymous Submission}

% The abstract must come before \maketitle (acmart).
% DRAFT: revise after Section 6 is final.
\begin{abstract}
Almost every machine-learning project now starts by choosing a pre-trained base model to adapt. On Hugging Face these
choices are extremely concentrated: ten models receive \topTenReal\% of first adoptions. We ask whether the
concentration reflects popularity, general quality, or fit to what developers need. We model each of \nFirst{} first
adoptions as a discrete choice over the whole market of earlier models (median \marketMedian{} alternatives). The model
is tractable because we sample competitors with probability proportional to degree plus one; McFadden's correction for
sampled alternatives keeps the estimator consistent, and the popular incumbents enter every choice set. A Shapley
decomposition of the likelihood gain attributes \shPopularity\% to popularity signals, \shMatch\% to
developer-specific demand match, \shInertia\% to inertia and almost nothing to reported benchmark quality. Fixed effects
for the 500 most adopted models bound the role of quality: even if every fixed attribute of these models were quality,
fit would account for at most \feInQual\% of the gain in sample and \feHoldQual\% out of sample, and external benchmark
scores explain only \extQ\% of the fixed effects. Among identical copies of the same weights, developers follow recent
adoptions, not total adoptions. A network grown forward from the estimated utility lies between a network grown from
popularity alone and one grown from fit alone, and is far closer to the former.
\end{abstract}

\begin{document}

\begin{titlepage}
\maketitle
\vspace{1cm}
\setcounter{tocdepth}{2}
\tableofcontents
\end{titlepage}

\section{Introduction}
\label{sec:intro}
\todo{Write after Section 6 is final. Planned structure for an EC audience:
(1) Setting: base-model choice is a technology-adoption decision in a market with hundreds of thousands of
products, little reliable quality information and highly visible popularity signals (downloads, likes, derivative
counts). The network of these choices on Hugging Face is extremely concentrated (\topTenReal\% of first adoptions go
to ten parents).
(2) Question: is the concentration driven by popularity (herding, cumulative advantage), by a few products of
exceptional quality, or by products that fit common needs? The answer matters for platform design (what a hub should
rank by) and for whether the AI supply chain concentrates efficiently.
(3) Approach: demand estimation over the entire market; sampling of alternatives with probability proportional to
degree plus one, with a consistency result (Proposition~\ref{prop:pps}) and an efficiency argument; Shapley
decomposition into popularity, quality, demand match and inertia; a fixed-effect bound on quality; counterfactual
growth of the network.
(4) Findings: popularity \shPopularity\%; developer-specific fit about \jointOwn\%; quality bounded and mostly not
benchmark quality; momentum among identical copies; growth result (Section 6).
(5) Contributions: (i) a sampling design and estimator for discrete choice over very large, heavy-tailed markets;
(ii) the decomposition and bound; (iii) evidence on popularity versus fit in a large technology-adoption market;
(iv) counterfactual networks.}

\section{Related Work}
\label{sec:related}

\paragraph{Popularity, social learning and herding.}
When quality is hard to observe, adopters rationally follow earlier adopters
\citep{banerjee1992simple,bikhchandani1992theory}, and social influence makes success both more unequal and less
predictable from quality \citep{salganik2006experimental}. Cumulative advantage in science has the same structure
\citep{merton1968matthew}. Field experiments show that displaying popularity information shifts choices
\citep{cai2009observational,tucker2011popularity}. In observational data, the central difficulty is that popularity
is itself a signal of quality. We address it in three ways: by bounding quality with product-specific fixed effects,
the largest share any unobserved fixed attribute could take; by measuring how much of that bound external benchmark
scores account for; and by comparing identical copies of the same weights, where quality is held fixed by
construction.

\paragraph{Demand estimation with very large choice sets.}
Conditional logit \citep{mcfadden1974conditional,train2009discrete} remains consistent when the likelihood is
evaluated on a sample of the alternatives, provided the sampling probabilities enter as an offset
\citep{mcfadden1978modelling}; related results cover semiparametric models \citep{fox2007semiparametric} and
multivariate extreme-value models \citep{guevara2013sampling}. These results allow any positive sampling
probabilities. Our setting, with hundreds of thousands of alternatives whose demand is concentrated on a few dozen,
makes the choice of design consequential. We sample with probability proportional to degree plus one, so that the
incumbents developers actually weigh are present in every choice set, and we show that a uniform design changes the
estimated weights and does not only widen the intervals.

\paragraph{Network formation as discrete choice.}
Preferential attachment explains heavy-tailed degree distributions by new nodes linking to existing nodes in
proportion to their degree \citep{price1976general,barabasi1999emergence}. Fitness models add a node-specific
attractiveness that lets late entrants overtake incumbents \citep{bianconi2001competition,wang2013quantifying}, and
empirical work measures the attachment kernel from the timing of new links
\citep{newman2001clustering,jeong2003measuring}. In economics, \citet{jackson2007meeting} show how a mix of random
and network-based search produces such degree distributions, and econometric models of network formation estimate
link preferences from observed networks \citep{graham2017econometric,depaula2020econometric}.
\citet{overgoor2019choosing} show that the physics mechanisms are conditional-logit models in which each new edge is a
choice among the existing nodes, so that degree, fitness, homophily and triadic closure can be estimated jointly and
compared by likelihood; importance sampling extends this to large graphs \citep{overgoor2020scaling}. We build on this
framework. We compare mechanisms with a Shapley decomposition of the likelihood gain instead of a sequence of nested
models, and we use the estimated utility to grow counterfactual networks.

\paragraph{AI model ecosystems.}
Foundation models are adapted rather than trained from scratch \citep{bommasani2021opportunities}, which makes the
choice of a base model a consequential and recurring decision. Descriptive studies of the Hugging Face Hub document
its growth, its family trees and the concentration of downstream use on a few model families
\citep{osborne2024ai,laufer2025anatomy}. In machine learning, base-model choice is usually framed as a prediction
problem: transferability estimators such as LEEP and LogME rank pre-trained models by how well they are expected to
fine-tune on a target task \citep{nguyen2020leep,you2021logme}. Our question is complementary and behavioral. We ask
not which model a developer \emph{should} choose but which signals the choices developers actually make respond to,
and what kind of market those choices grow.

\section{Data and Setting}
\label{sec:data}

\paragraph{The fine-tuning network.}
We use a snapshot of the Hugging Face Hub with metadata and model cards for \nModels{} models by \nAuthorsAll{}
accounts.\todo{Cite the snapshot source and its date.} A model card may declare a base model; the declaration and
its type define a directed edge from the parent (base) to the child. Of the \nEdges{} edges, \nAdapter{} are
adapters and \nFinetune{} full fine-tunes; the remainder are quantized copies and merges. We treat fine-tunes and
adapters (and the few untyped declarations) as \emph{base-model choices}: a developer picked one existing model to
build on. Quantized copies and merges are not choices of this kind; they enter only as attributes of the parent
(e.g., how many quantized versions a candidate has). Creation timestamps are reliable from \tStart{} onward; older
models are kept as part of the initial network, and every child created between \tStart{} and \tEnd{} contributes
one choice per declared parent, \nChoices{} choices in total.

\paragraph{From uploads to adoption decisions.}
Many uploads are repeated variants of one decision: the same author fine-tunes the same base model several times
within hours (checkpoints, hyper-parameter sweeps). We merge consecutive uploads by one author on one parent that
are less than 24 hours apart into a \emph{decision} (\nDecisions{} decisions; 6- and 72-hour windows in
\cref{app:repeat}). Our main sample is the first decision of each author--parent pair, a \emph{first adoption}:
\nFirst{} first adoptions by \nFirstAuthors{} authors of \nFirstParents{} distinct parents. First adoptions are the
decisions in which a developer selects a base model rather than continuing to use one; repeated use is a different
behavior, which we analyze separately in \cref{app:repeat}. The \emph{market} for a first adoption is every model of
the same modality (text, image generation, speech, vision, multimodal) created before the child, minus the parents
the author has already used. Markets are large: the median first adoption faces \marketMedian{} alternatives, and
\shareText\% of first adoptions are text models.

\paragraph{A concentrated network.}
Adoption is extremely concentrated. The single most adopted parent receives \topOneReal\% of all first adoptions,
the ten most adopted receive \topTenReal\%, and the effective number of parents, $\exp(H)$ for the entropy $H$ of
the adoption shares, is \effReal{} out of \nFirstParents{} parents ever adopted. Concentration is what popularity
would produce, but also what a few models of exceptional general quality, or a few models that fit the most common
needs, would produce. Separating these explanations is the purpose of the choice model.

\paragraph{What a developer can see: four groups of signals.}
For every (choice, candidate) pair we measure signals available to the developer at the time of the choice, from the
network state and metadata strictly before the child's creation. We organize them into four groups
(\cref{app:data} lists every feature):
\begin{itemize}
\item \textbf{Popularity}: signals that are the same for every developer and do not depend on what the developer
needs: degree $d$ and $\1[d>0]$, adoptions in the last 30 days, age, a major-organization publisher, a legacy
model (created before March 2022), the library ecosystem, the number of quantized versions, re-distributor accounts that republish
other organizations' weights, co-use by other developers who trained on the same datasets or used the same niche
parents, and whether size, card text or benchmark scores are available at all (documentation completeness is
visibility, not demand).
\item \textbf{General quality}: model-index benchmark scores reported in the card, standardized within benchmark
(a z-score). Hub fixed effects, which absorb any unobserved quality, are used only to bound this group
(\cref{sec:bounds}).
\item \textbf{Demand match}: how well the candidate fits this developer's need. Need is read only from sources the
child does not copy from its parent: the author's earlier parents (task, languages, size; the candidate itself left
out), the languages and tasks common among other models trained on the child's declared datasets, and the TF-IDF
similarity between the child's and the candidate's model cards.
\item \textbf{Inertia}: the candidate is the author's own model, or belongs to a model family the author used before.
\end{itemize}

\paragraph{Guarding against copied metadata.}
A child's card is partly inherited from its parent, so naive similarity between child and candidate would favor the
true parent mechanically. We test each field for inheritance by comparing the child with its parent and with a
random parent from the same month and modality (Cohen's $\kappa$, \cref{app:data}). Parameter count ($\kappa=0.84$)
and license ($\kappa=0.74$) are mostly copied and are not used as demand; datasets ($\kappa=0.12$) are mostly the
child's own. For dataset-based demand we additionally drop datasets the child shares with its chosen parent, compute
each dataset's language and task profile only from models created before the choice, and remove the candidate's own
children from the profile used for that candidate. For card text we drop template paragraphs (paragraphs in
$\ge$20 cards, including near-duplicates found by MinHash), paragraphs copied from an earlier card, and organization
and family names, so that a shared model name cannot produce similarity.

\paragraph{Training and holdout.}
We hold out the last 10\% of choices in time (children created on or after \tHold) for out-of-sample evaluation;
\nTrain{} first adoptions are used for estimation. The fixed effects of \cref{sec:bounds} are assigned to the 500
most adopted parents as of the holdout cut, so the set does not depend on future adoptions.

\section{Model and Estimation}
\label{sec:method}

\subsection{Base-model selection as a discrete choice}
Following \citet{overgoor2019choosing}, each first adoption $k$ is a choice of one parent $j_k$ from the market
$\M_k$. Developer $k$ values candidate $j$ at $U_{kj} = \theta^\top x_{kj} + \varepsilon_{kj}$ with i.i.d.\ Gumbel
errors, which gives the conditional logit \citep{mcfadden1974conditional}
\begin{equation}
  P(j \mid \M_k) = \frac{\exp(\theta^\top x_{kj})}{\sum_{l \in \M_k} \exp(\theta^\top x_{kl})}.
  \label{eq:logit}
\end{equation}
The features $x_{kj}$ are the four groups of \cref{sec:data}; they vary over candidates and over time and, for
demand match and inertia, over developers. With $\theta = 0$ the model is \emph{uniform attachment}, every
alive model equally likely; preferential attachment is the special case in which only $\log d$ and $\1[d>0]$
enter. We interpret $\theta$ descriptively, as the weights that best summarize which signals choices co-move with,
not as causal effects (\cref{sec:discussion}).

\subsection{Sampling the competitors that matter}
\label{sec:sampling}
Evaluating \eqref{eq:logit} over a median of \marketMedian{} alternatives for \nFirst{} choices is infeasible, so we
sample a candidate set $\D_k \subset \M_k$. Each alternative enters independently (Poisson sampling) with
\begin{equation}
  \pi_{kl} = \min\{1,\; c_k (d_{kl} + 1)\}, \qquad \textstyle\sum_{l \in \M_k} \pi_{kl} = S,
  \label{eq:pps}
\end{equation}
where $d_{kl}$ is the degree of $l$ just before choice $k$ and $c_k$ is solved in closed form from the degree
histogram of the market; the chosen parent is always included. We use $S=300$; on average 8.6 candidates per choice
have $\pi=1$, so the incumbents a developer is most likely to weigh are present in every set, and the rest of the
set is a probability-proportional-to-size sample of the long tail.

\begin{proposition}[Sampling correction]
\label{prop:pps}
Under \eqref{eq:logit} and the sampling scheme \eqref{eq:pps},
\begin{equation}
  P(j \mid \D_k) = \frac{\exp(\theta^\top x_{kj} - \log \pi_{kj})}{\sum_{l \in \D_k} \exp(\theta^\top x_{kl} - \log
  \pi_{kl})},
  \label{eq:corrected}
\end{equation}
and maximizing $\sum_k \log P(j_k \mid \D_k)$ gives a consistent estimator of $\theta$.
\end{proposition}
The proof (\cref{app:sampling}) applies McFadden's positive-conditioning argument \citep{mcfadden1978modelling}:
given that $j$ was chosen, the probability of drawing $\D$ is $\prod_{l \in \D \setminus j} \pi_l
\prod_{l \notin \D}(1-\pi_l)$, so the ratio for two members of $\D$ is $\pi_l / \pi_j$. The offset $-\log \pi$
makes uniform attachment, $\theta = 0$, the natural baseline: it assigns every sampled candidate the weight
$1/\pi$, i.e.\ every model in the full market the same probability.

Why sample proportionally to degree? Consistency holds for any $\pi > 0$, but efficiency does not. A uniform sample
of 300 from a market of $3 \times 10^5$ almost never contains the dozen models that receive most adoptions, so each
choice compares the chosen parent mostly with never-used models, and the likelihood learns little about how
developers trade off one popular model against another. \Cref{sec:robustq} shows that this is not only a matter of
precision: with uniform sets the estimated weights change.

\subsection{Estimation and inference}
The training data have 35.2 million (choice, candidate) rows. We group choices into buckets of similar set size,
keep the 26 features on the GPU in single precision and accumulate the log-likelihood, gradient and Hessian in double
precision; L-BFGS with an analytic Hessian fits the full model in under a minute on one A100. Standard errors are
clustered by author \citep{cameron2015practitioner}: \nFirstAuthors{} authors make the \nFirst{} first adoptions, and
the choices of one author are strongly dependent (\cref{app:validation}). For every derived quantity, including the
Shapley shares below, we use a multiplier (Bayesian) bootstrap over authors \citep{rubin1981bayesian}: each
replicate draws a weight $w_a \sim \mathrm{Exp}(1)$ per author, reweights the author's choices and refits from the
point estimate.

\subsection{Decomposing the likelihood gain}
\label{sec:shapley}
The four groups are correlated (popular models also tend to be well documented and to match common tasks), so the
change in fit from adding a group depends on which groups are already in the model. We therefore attribute the
log-likelihood gain over uniform attachment, $G(\mathcal{S}) = \ell(\hat\theta_{\mathcal{S}}) - \ell(0)$ for the
model with the groups in $\mathcal{S}$, by the Shapley value \citep{shapley1953value}
\begin{equation}
  \phi_g = \sum_{\mathcal{S} \subseteq \mathcal{G} \setminus \{g\}}
  \frac{|\mathcal{S}|!\,(|\mathcal{G}|-|\mathcal{S}|-1)!}{|\mathcal{G}|!}
  \big[G(\mathcal{S} \cup \{g\}) - G(\mathcal{S})\big],
  \label{eq:shapley}
\end{equation}
and report shares $\phi_g / G(\mathcal{G})$. With four groups this needs $2^4 = 16$ fits, each warm-started from
a subset one group smaller. We compute the shares both in-sample and on the holdout (mean per-choice
log-likelihood), where a group can only gain by predicting future choices.

\subsection{Bounding general quality}
\label{sec:bounds}
Benchmark scores cover few models, so a small quality share could reflect missing measurement rather than
indifference to quality. We bound the quality share from above by adding an intercept $\phi_j$ for each of the
500 most adopted parents (assigned ex ante, \cref{sec:data}). The fixed effects absorb every time-invariant
attribute of a hub: unmeasured quality, but also brand, tutorials and default status. Entering the fixed effects as
a fifth Shapley player and counting their entire share as quality gives the largest share that any fixed hub
attribute could contribute. To see how much of this bound is quality, we regress $\hat\phi_j$ on external quality
measures (Open LLM Leaderboard v2 average, arena ratings; controls for parameter count and re-distributor accounts)
and correct the $R^2$ for estimation noise using the bootstrap variance of $\hat\phi_j$; the noise-corrected share
of $\mathrm{Var}(\phi)$ explained by quality splits the fixed-effect share into a quality and a popularity part.

\subsection{Does the decomposition recover known mechanisms?}
\label{sec:validation}
Two checks support the procedure (details in \cref{app:validation}). First, on synthetic preferential-attachment
networks the corrected likelihood recovers the attachment exponent under all three sampling designs. Second, we
replace every observed choice with a draw from a known utility on the same candidate sets and re-run the full
decomposition. When choices are drawn from the popularity terms alone, the decomposition attributes
\recSfPop\% of the gain to popularity; when drawn from quality and demand match alone, it attributes \recQfFit\% to
fit; when drawn from the full estimated model, it returns \recFullPop\%, \recFullMatch\% and \recFullInertia\% for
popularity, demand match and inertia, the shares estimated on the real choices. The decomposition can tell the
mechanisms apart, and it does not manufacture a popularity share where none was generated.

\section{What Drives Base-Model Choice}
\label{sec:results}

\begin{table}[t]
\centering
\caption{Shapley decomposition of the log-likelihood gain over uniform attachment (main specification: $(d{+}1)$ PPS candidate sets, first adoptions). In-sample shares with 95\% multiplier-bootstrap intervals (50 replicates, authors resampled); holdout = mean per-choice log-likelihood on the last 10\% of choices in time.}
\label{tab:shapley}
\small
\begin{tabular}{lrcr}
\toprule
Group & Share (\%) & 95\% CI & Holdout (\%) \\
\midrule
Popularity & 91.6 & [91.1, 92.0] & 91.8 \\
General quality & 0.006 & [0.003, 0.010] & 0.001 \\
Demand match & 5.2 & [5.0, 5.4] & 4.6 \\
Inertia & 3.2 & [2.7, 3.6] & 3.6 \\
\midrule
\multicolumn{4}{l}{Holdout gain over uniform: 6.62 nats/choice; top-1 / top-10 accuracy 14.2\% / 39.8\%} \\
\bottomrule
\end{tabular}
\end{table}


\subsection{Popularity explains almost all of what can be explained}
\Cref{tab:shapley} reports the decomposition for the main specification. Popularity accounts for
\shPopularity\% of the log-likelihood gain over uniform attachment (95\% CI \shPopularityLo--\shPopularityHi),
demand match for \shMatch\% (\shMatchLo--\shMatchHi), inertia for \shInertia\% (\shInertiaLo--\shInertiaHi) and
general quality for less than 0.01\%. The holdout shares are nearly identical (\shHoldPopularity\%, \shHoldMatch\%,
\shHoldInertia\%), so the ranking is not an artifact of in-sample fit. Every group is jointly significant (likelihood
ratio for dropping popularity \lrPop, demand match \lrMatch, inertia \lrInertia, quality \lrQual{} on 1 d.f.); the
question is not whether developers respond to fit but how much of their choices fit accounts for. Out of sample the
model gains \holdGain{} nats per choice over uniform attachment and ranks the chosen parent first among the
$\approx$300 sampled candidates in \holdTopOne\% of holdout choices (top-10: \holdTopTen\%).

\paragraph{Developers value fit, but fit rarely separates candidates.}
The small share of demand match is not because developers ignore it. In utility terms (\cref{tab:coef}), a
candidate that matches the task of the author's earlier parents is worth as much as a \mrsTask-fold increase in
$1+d$ (95\% CI \mrsTaskLo--\mrsTaskHi); one standard deviation of recent adoptions is worth an \mrsRec-fold increase,
a major-organization publisher \mrsOrg-fold. Per unit, fit is valued highly. It explains little of the variation in
choices because the two kinds of signal vary very differently across candidates: degree spans five orders of
magnitude within a candidate set, whereas the match signals are coarse, and many candidates fit equally well. Among
first adoptions with a demand signal, on average \tieMean{} sampled candidates tie at the highest match index, and
what separates them is popularity. Fit narrows the field; popularity picks the winner.

\paragraph{Quality.}
Model-index benchmark scores are reported for well under 1\% of candidates, and their coefficient is small and, if
anything, negative. This could mean that developers do not respond to quality, or only that quality is rarely
measured in the card. The next subsection bounds the second possibility.

\begin{table}[t]
\centering
\caption{How large could fit be? Fit = general quality + demand match. Hub FE: 500 hub-specific intercepts (top-500 by degree at the holdout cut). External quality explains 7.3\% of the FE variance among the 118 hubs covered by the Open LLM Leaderboard v2 (noise-corrected, with redistributor and size controls).}
\label{tab:bounds}
\small
\begin{tabular}{lrr}
\toprule
Specification & Fit (\%) & Popularity (\%) \\
\midrule
No FE; quality = benchmark z-score & 5.2 & 91.6 \\
Hub FE counted as quality (in-sample) & 46.0 & 51.0 \\
Hub FE counted as quality (holdout) & 36.2 & 60.3 \\
Externally explained part of FE (in-sample) & 7.0 & 89.9 \\
Externally explained part of FE (holdout) & 6.3 & 90.2 \\
\bottomrule
\end{tabular}
\end{table}


\subsection{How large could fit be?}
\paragraph{An upper bound from hub fixed effects.}
Adding fixed effects for the 500 most adopted parents, and counting them entirely as quality, is the most generous
reading of quality: any fixed attribute of a hub, measured or not, would be credited to it. Under this reading fit
would account for \feInQual\% of the in-sample gain and \feHoldQual\% of the holdout gain (\cref{tab:bounds}). Two
observations show that the bound is far from tight. First, the fixed effects add almost nothing out of sample:
the holdout gain moves from \holdGainNoFE{} to \holdGainFE{} nats per choice. In-sample, the fixed effects take
over credit from the popularity features they are collinear with (the popularity share falls from \shPopularity\%
to 51.0\%), but they predict future choices no better than the popularity signals already do. Second, external
quality explains little of them: among the \extHubs{} hubs covered by the Open LLM Leaderboard v2, the
noise-corrected share of fixed-effect variance explained by the leaderboard score (with size and re-distributor
controls) is \extQ\%. Splitting the fixed-effect share accordingly, fit accounts for \extFitIn\% in-sample and
\extFitHold\% on the holdout. Whatever makes a hub attractive beyond its observed popularity, it is mostly not what
benchmarks measure.\todo{Caveat: leaderboard coverage is 118 of 500 hubs; discuss representativeness.}

\paragraph{Is the demand-match share developer-specific demand?}
Some of what looks like match could be common trends: if everyone is fine-tuning instruction-following text models
in a given month, any developer's history will ``match'' the models that are popular that month. We replace each
developer's demand with that of a donor, another author's first adoption in the same month and modality, and
recompute the match features. On its own the donor's demand still accounts for \placeboAlone\% of the gain, so common
trends are real. But when own and donor match enter the same model, own match keeps \jointOwn\% of the gain and the
donor \jointDonor\%, and the own-match coefficients barely move (task history 1.93 vs.\ 1.90). The donor's demand
competes with popularity, not with the developer's own demand: about \jointOwn\% of the gain is fit to a
developer-specific need.

\paragraph{Popularity without a quality difference.}
The cleanest evidence that popularity signals act on their own comes from \idN{} first adoptions in which two copies
of the same weights (identical weight hashes, published by different accounts) are both in the candidate set and
one is chosen. Quality is identical by construction. The chosen copy is the one with more recent adoptions
(coefficient \idRec, SE \idRecSe), not the one with more adoptions overall (\idDeg, SE \idDegSe; the chosen copy has
the higher degree in \idHigher\% of cases): among exact substitutes, developers follow current momentum.

\begin{table*}[t]
\centering
\caption{The decomposition under three candidate-set designs. Same choices, features and estimator; only the inclusion probabilities $\pi$ differ. Coefficients with bootstrap SEs. Under uniform sampling almost every sampled competitor is a never-used model, so top-1 accuracy is high but the data contain few comparisons among popular models.}
\label{tab:q}
\footnotesize
\setlength{\tabcolsep}{4pt}
\begin{tabular}{lccc}
\toprule
 & $(d{+}1)$ PPS & Stratified & Uniform \\
\midrule
Popularity share (\%) & 91.6 [91.1, 92.0] & 92.2 [92.0, 92.5] & 88.3 [87.8, 88.6] \\
General quality share (\%) & 0.006 [0.003, 0.010] & 0.006 [0.003, 0.010] & 0.009 [0.005, 0.015] \\
Demand match share (\%) & 5.2 [5.0, 5.4] & 5.6 [5.4, 5.9] & 8.4 [8.1, 8.7] \\
Inertia share (\%) & 3.2 [2.7, 3.6] & 2.1 [1.9, 2.4] & 3.4 [3.1, 3.8] \\
\midrule
$\log d$ & 0.34 (0.01) & 0.34 (0.01) & 0.44 (0.03) \\
$\mathbf{1}[d>0]$ & 4.64 (0.09) & 4.71 (0.09) & 3.25 (0.10) \\
major org. & 0.65 (0.03) & 0.68 (0.03) & 1.48 (0.05) \\
language history & 0.25 (0.03) & 0.24 (0.03) & 0.95 (0.06) \\
task history & 1.90 (0.04) & 1.88 (0.04) & 2.38 (0.08) \\
card-text sim. & 1.69 (1.64) & 4.36 (0.90) & 0.74 (2.66) \\
\midrule
Holdout top-1 acc. (\%) & 14.2 & 13.5 & 93.1 \\
Bootstrap replicates (shares / coef.) & 50 / 200 & 20 / 80 & 20 / 80 \\
\bottomrule
\end{tabular}
\end{table*}


\subsection{Robustness to the candidate-set design}
\label{sec:robustq}
\Cref{tab:q} re-estimates the model with the same choices and features under three inclusion schemes: the main
$(d+1)$ PPS design, a stratified design (all top-500 hubs plus random positive-degree and zero-degree models) and
uniform sampling. PPS and stratified sets give the same answer: shares within one percentage point and coefficients
within one standard error. Both contain the popular competitors, and the $-\log\pi$ correction handles the rest.

Uniform sets give different weights. Almost every sampled competitor is a never-used model, so the chosen parent is
easy to find (top-1 holdout accuracy \topOneUnif\%), but the choice is then mostly a comparison between ``known'' and
``unknown'' models. On this margin the weights on organization brand, language history and the library ecosystem
roughly double, $\1[d>0]$ falls from 4.64 to 3.25, and the match share rises to 8.4\%. Since all three designs
recover the true parameters on synthetic data (\cref{sec:validation}), the difference is not an estimation
artifact: a single logit cannot describe both the known-versus-unknown margin and the choice among popular models
with the same weights. We take the PPS design as the main specification because it represents the comparisons
developers actually make, among the incumbents of their modality.

\paragraph{Other checks.}
Demand match can only register where a demand signal exists; \sigShare\% of first adoptions have one (author
history, dataset profile or card text). On these choices demand match accounts for \sigMatch\% of the gain
(\sigMatchRw\% after reweighting them to the composition of all choices by author experience and organization), and
popularity for \sigPop\%. Full-data estimates agree with an earlier analysis of a 10\% sample: \valIn{} of \valN{}
coefficients lie in the 10\% confidence interval, and the exceptions are features redefined since
(\cref{app:validation}).

\section{Growing Counterfactual Networks}
\label{sec:growth}
\todo{Results are in (results/tables/v2/v2\_grow.md, v2\_grow\_recovery.md); the over-concentration diagnosis is
pending. Planned content:
(i) The growth simulator: the real event sequence (who adopts, when, in which modality) is kept; parents are redrawn
from PPS sets built on simulated degrees; all 26 features are recomputed against the simulated state (exact replay
of the estimation features as a check).
(ii) Networks: SF (popularity only), QF (quality and match only), EST (estimated utility), each with and without
hub fixed effects, 5 seeds. Effective number of parents: SF 439, EST 509, real 837, QF 114,605. With FE: SF 400,
EST 429. EST lies between SF and QF and much closer to SF.
(iii) Share recovery on the grown networks: SF 98.5\% popularity; QF 92.5\% demand match; EST 94.3\% / 4.7\% /
0.9\% (popularity / match / inertia).
(iv) Feedback over-concentrates relative to the real network (top-10 share 0.31--0.39 against 0.22). Candidate
explanation: exogenous entry of new model generations.
(v) Counterfactual that scales the popularity weights.}

\section{Discussion and Limitations}
\label{sec:discussion}
\todo{Write last. Points:
(1) The estimates are descriptive, not causal. Popularity signals may proxy for information developers have that we
do not observe (word of mouth, documentation outside the card). The identical-copy comparison is the closest we get
to holding quality fixed.
(2) Social learning or herding? Following popularity is rational when quality is hard to observe; the
identical-copy result (recent adoptions matter, total adoptions do not, quality held fixed) separates momentum from
learning about quality.
(3) Platform design: ranking by downloads and likes reinforces the dominant mechanism. Surfacing fit-relevant
information (tasks, languages, datasets) could matter more than surfacing quality scores. Link to the counterfactual
of Section 6.
(4) Concentration of the AI supply chain: most downstream models inherit from a few bases, which matters for
licensing, safety and market power.
(5) Demand is measured only where a signal exists (\sigShare\% of first adoptions).
(6) Single snapshot: deleted models and edited cards are not observed.
(7) The quality bound uses leaderboard coverage of 118 of 500 hubs.}


\bibliographystyle{ACM-Reference-Format}
\bibliography{refs}

\appendix
\section{Data and Features}
\label{app:data}

\paragraph{Network construction.}
Nodes are models; an edge $p \to c$ is a base-model declaration in the card of $c$. We drop edges whose child was
created before its parent and equal-timestamp edges outside the pre-2022 backfill, and keep deleted parents with at
least 100 children as placeholder nodes. Creation times up to March 2, 2022 are backfilled in the source data; models
created up to that date are \emph{legacy} models, part of the initial network, and are given that date as their
creation time. A choice is an edge of type fine-tune, adapter or untyped whose child is created after that date and
whose parent has the child's modality; modality is taken from the pipeline tag, else from the parent, else from the
library.

\paragraph{Features.}
All features of choice $k$ use only models and edges created strictly before the child of $k$. ``History'' is the
set of parents the author adopted before $k$, with the candidate itself left out.

\begin{center}
\scriptsize
\begin{tabular}{llp{0.62\textwidth}}
\toprule
Group & Feature & Definition \\
\midrule
Popularity & \texttt{log\_deg}, \texttt{has\_deg} & $\log \max(d,1)$ and $\1[d>0]$, $d$ = number of earlier fine-tune/adapter children \\
 & \texttt{lrec30} & $\log(1+{}$children in the last 30 days$)$ \\
 & \texttt{log\_age} & $\log(1+{}$age in days$)$ \\
 & \texttt{major\_org} & published by one of a fixed list of large AI organizations \\
 & \texttt{legacy} & created before March 2, 2022 \\
 & \texttt{lquant} & $\log(1+{}$quantized children$)$ \\
 & \texttt{lib\_*} & library: transformers, peft, diffusers, sentence-transformers \\
 & \texttt{redistributor} & published by an account whose models are mostly copies of other organizations' models ($\ge$20 copies, $\ge$50\% of its models) \\
 & \texttt{ds\_couse} & $\log(1+{}$earlier models by other authors, trained on one of the child's datasets, that chose the candidate$)$ \\
 & \texttt{lco2} & $\log(1+{}$distinct other authors who share a niche parent with the author and used the candidate$)$; niche = outside the month's top-500 \\
 & \texttt{size\_na}, \texttt{txt\_na}, \texttt{bench\_na} & candidate size (or the author's size history), usable card text, or benchmark score missing \\
\midrule
Quality & \texttt{bench\_z} & model-index benchmark score in the candidate's card, z-scored within benchmark (0 if missing) \\
\midrule
Demand match & \texttt{task\_hist} & candidate's pipeline task appears in the author's history \\
 & \texttt{lang\_hist} & candidate shares a language with the author's history \\
 & \texttt{size\_dist} & $|\log \text{params}(j) - \text{median} \log \text{params}(\text{history})|$ \\
 & \texttt{ds\_lang}, \texttt{ds\_task} & candidate's language / task is common ($\ge$30\%) among other models trained on the child's datasets (profiles as of the choice; datasets shared with the chosen parent dropped; the candidate's own children removed) \\
 & \texttt{txt\_sim} & TF-IDF cosine between child and candidate cards after template, copy and name removal; 0 unless both have $\ge$5 tokens \\
\midrule
Inertia & \texttt{same\_author} & candidate published by the choosing author \\
 & \texttt{author\_family} & candidate in a model family (root) the author used before \\
\bottomrule
\end{tabular}
\end{center}

\paragraph{Inheritance check.}
For each metadata field we compare the child with its parent and with a random parent of another choice in the same
modality and month (\cref{tab:inherit}). Fields with high excess agreement are inherited and cannot measure the
developer's demand; demand match therefore uses the author's history, dataset profiles and cleaned card text only.

\begin{table}[h]
\centering
\caption{Inheritance of card metadata, first adoptions: agreement between child and parent, between child and a
random parent (chance), and Cohen's $\kappa$.}
\label{tab:inherit}
\small
\begin{tabular}{lrrrr}
\toprule
Field & Both declare & Identical & Chance & $\kappa$ \\
\midrule
Parameter count & 0.25 & 0.86 & 0.11 & 0.84 \\
License & 0.73 & 0.83 & 0.35 & 0.74 \\
Languages & 0.23 & 0.75 & 0.55 & 0.44 \\
Pipeline task & 0.57 & 0.69 & 0.52 & 0.34 \\
Datasets & 0.08 & 0.13 & 0.01 & 0.12 \\
Library & 0.90 & 0.68 & 0.65 & 0.07 \\
\bottomrule
\end{tabular}
\end{table}

\paragraph{Card text.}
We clean 1,062,460 cards. A paragraph is a template if it (or a MinHash near-duplicate, 5-word shingles, Jaccard
$\ge 0.8$) appears in at least 20 cards; it is copied if it appears in a card created no later than the child's.
Organization and family names (the first token of a repository name) and tokens with digits are removed. After
cleaning, 18\% of child cards keep at least 5 tokens, the threshold for \texttt{txt\_sim}.

\section{Sampling Correction}
\label{app:sampling}

\paragraph{Proof of \cref{prop:pps}.}
Fix choice $k$ and drop the index. Every $l \neq j$ enters $\D$ independently with probability $\pi_l$, and the
chosen $j$ is always included. Given that $j$ was chosen, the probability of a set $\D \ni j$ is
$q(\D \mid j) = \prod_{l \in \D \setminus \{j\}} \pi_l \prod_{l \notin \D} (1 - \pi_l)$. Hence for $j, l \in \D$,
$q(\D \mid j) / q(\D \mid l) = \pi_l / \pi_j$. By Bayes' rule
\begin{equation*}
P(j \mid \D) = \frac{P(j)\, q(\D \mid j)}{\sum_{l \in \D} P(l)\, q(\D \mid l)}
= \frac{e^{\theta^\top x_j}\, q(\D \mid j)}{\sum_{l \in \D} e^{\theta^\top x_l}\, q(\D \mid l)}
= \frac{e^{\theta^\top x_j - \log \pi_j}}{\sum_{l \in \D} e^{\theta^\top x_l - \log \pi_l}},
\end{equation*}
which is \eqref{eq:corrected}. The conditional likelihood is a correctly specified likelihood for the observed
$(j_k, \D_k)$, and $q(\D \mid j) > 0$ for every $j \in \D$ (positive conditioning), so the usual M-estimation
argument gives consistency \citep{mcfadden1978modelling,train2009discrete}. \hfill$\square$

\paragraph{Solving for $c_k$.}
Let $n_0$ be the number of degree-0 models in the market and $n_\delta$ the number with degree $\delta > 0$. Then
$\sum_l \pi_l = n_0 \min(1,c) + \sum_\delta n_\delta \min\{1, c(\delta + 1)\}$ is piecewise linear and increasing in
$c$. Sorting the positive degrees in decreasing order, the top $m$ degree classes are saturated ($\pi=1$) and
$c = (S - \sum_{i \le m} n_{\delta_i}) / (n_0 + \sum_{i > m} n_{\delta_i} (\delta_i + 1))$ for the unique $m$ at which
this $c$ saturates exactly the top $m$ classes. The histogram is maintained incrementally, so $c_k$ costs one pass
over the distinct degrees. Degree-0 models are drawn as a binomial count followed by rejection sampling.

\paragraph{Author history and plan A.}
The candidate sets also contain the author's earlier parents with $\pi = 1$ (used for repeated choices,
\cref{app:repeat}). For first adoptions these rows are removed, consistent with defining the market as the models the
author has not used.

\section{Validation}
\label{app:validation}

\paragraph{Synthetic preferential attachment.}
We grow synthetic networks of 8,000 nodes in which each new node attaches to an existing node with weight
$d^{\alpha}$ if its degree $d$ is positive and $e^{\beta_0}$ otherwise ($\alpha = 1$, $\beta_0 = -1$), replay them with the same code as the real data and estimate
$(\1[d>0], \log d)$ under each of the three sampling designs. In all three, $\hat\alpha$ lies within three standard
errors of 1. The empirical inclusion frequency of each node under PPS matches $\pi$, and the expected set size
matches $S$.

\paragraph{Recovering simulated mechanisms.}
For \cref{sec:validation} we keep every real candidate set and replace the chosen parent by a draw with probability
proportional to $\exp(V_{kl} - \log \pi_{kl})$, where $V$ contains (a) only the popularity terms of the estimated
model, (b) only quality, demand match and the externally explained part of the hub fixed effects, or (c) the full
model with fixed effects. Re-running the four-group decomposition on these choices gives
\recSfPop\% popularity for (a), \recQfFit\% fit for (b), and \recFullPop\% / \recFullMatch\% / \recFullInertia\%
for (c). Because the candidate features are the real ones, the popularity and match signals keep their real
correlation; (a) and (b) therefore do not reach 100\%.

\paragraph{Full data versus a 10\% sample.}
An earlier version of the analysis used a 10\% random sample of choices. \Cref{tab:validate} compares it with the
full-data estimates under the stratified design (the design of the earlier analysis). \valIn{} of \valN{} coefficients
lie inside the 95\% interval of the 10\% sample. The exceptions are features changed in the rebuild: major
organization (the re-distributor flag now separates accounts that republish other organizations' weights), missing
card text (new text pipeline) and features based on dataset profiles, which are now computed only from earlier models.

\paragraph{Why standard errors shrink by less than $\sqrt{10}$.}
With independent choices the full-data standard errors would be $\sqrt{0.1} \approx 0.32$ of the 10\% ones; the
median ratio is \valSeRatio. Choices are clustered by author: the \nFirst{} first adoptions come from \nFirstAuthors{}
authors, and a 10\% sample of choices keeps 15\% of the authors. Under clustering the variance scales with
$\sum_a n_a^2 / N^2$, which bounds the ratio between 0.32 (independent choices within an author) and 0.89 (perfectly
dependent choices); the observed value implies strong within-author dependence, which is why every interval in the
paper is clustered or bootstrapped by author.

\begin{table}[htbp]
\centering
\caption{Full-data estimates (stratified design) against the earlier 10\% sample. Differences outside the 10\% interval concern features redefined in the rebuild.}
\label{tab:validate}
\scriptsize
\begin{tabular}{lccrr}
\toprule
Feature & 10\% sample & Full & In CI & SE ratio \\
\midrule
\texttt{author\_family} & -0.062 (0.085) & -0.065 (0.062) & yes & 0.73 \\
\texttt{bench\_na} & 0.176 (0.096) & -0.055 (0.048) & no & 0.50 \\
\texttt{bench\_z} & 0.015 (0.080) & -0.110 (0.025) & yes & 0.31 \\
\texttt{ds\_couse} & 1.150 (0.040) & 1.217 (0.020) & yes & 0.48 \\
\texttt{ds\_lang} & 0.459 (0.122) & 0.340 (0.051) & yes & 0.42 \\
\texttt{ds\_task} & 1.245 (0.149) & 0.970 (0.116) & yes & 0.78 \\
\texttt{has\_deg} & 4.687 (0.079) & 4.711 (0.082) & yes & 1.03 \\
\texttt{lang\_hist} & 0.360 (0.042) & 0.239 (0.027) & no & 0.64 \\
\texttt{lco2} & 0.676 (0.135) & 0.773 (0.095) & yes & 0.70 \\
\texttt{legacy} & 1.372 (0.052) & 1.379 (0.033) & yes & 0.64 \\
\texttt{lib\_diffusers} & 0.442 (0.269) & 0.555 (0.154) & yes & 0.57 \\
\texttt{lib\_peft} & -1.763 (0.305) & -1.674 (0.139) & yes & 0.46 \\
\texttt{lib\_sentence-transformers} & 0.873 (0.113) & 0.941 (0.047) & yes & 0.41 \\
\texttt{lib\_transformers} & 0.196 (0.082) & 0.202 (0.036) & yes & 0.45 \\
\texttt{log\_age} & -0.539 (0.021) & -0.540 (0.019) & yes & 0.90 \\
\texttt{log\_deg} & 0.347 (0.014) & 0.336 (0.011) & yes & 0.79 \\
\texttt{lquant} & 0.099 (0.013) & 0.070 (0.008) & no & 0.65 \\
\texttt{lrec30} & 0.631 (0.019) & 0.617 (0.015) & yes & 0.80 \\
\texttt{major\_org} & 0.300 (0.046) & 0.681 (0.031) & no & 0.68 \\
\texttt{same\_author} & 6.566 (0.449) & 6.319 (0.504) & yes & 1.12 \\
\texttt{size\_dist} & -0.173 (0.025) & -0.226 (0.017) & no & 0.66 \\
\texttt{size\_na} & -0.780 (0.062) & -0.824 (0.035) & yes & 0.56 \\
\texttt{task\_hist} & 1.909 (0.077) & 1.882 (0.044) & yes & 0.57 \\
\texttt{txt\_na} & -0.743 (0.074) & -0.368 (0.058) & no & 0.78 \\
\texttt{txt\_sim} & 4.559 (0.656) & 4.360 (0.757) & yes & 1.15 \\
\bottomrule
\end{tabular}
\end{table}


\section{Coefficients and Further Robustness}
\label{app:robust}

\begin{table}[h]
\centering
\caption{Conditional-logit coefficients, main specification (no FE). SE clustered by author; 95\% multiplier-bootstrap interval (200 replicates). Last column: the factor on $(1+d)$ that is worth the same utility as a one-SD (binary: 0$\to$1) change of the feature.}
\label{tab:coef}
\scriptsize
\begin{tabular}{lrrcr}
\toprule
Feature & Coef. & SE & 95\% CI & $(1{+}d)$ multiple \\
\midrule
\multicolumn{5}{l}{\textit{Popularity}} \\
\texttt{log\_deg} & 0.344 & 0.012 & [0.319, 0.364] & -- \\
\texttt{lrec30} & 0.605 & 0.015 & [0.582, 0.637] & 8.49 \\
\texttt{major\_org} & 0.650 & 0.031 & [0.591, 0.711] & 6.61 \\
\texttt{has\_deg} & 4.637 & 0.081 & [4.468, 4.769] & 7.05e+05 \\
\texttt{legacy} & 1.386 & 0.034 & [1.321, 1.444] & 55.9 \\
\texttt{log\_age} & -0.539 & 0.021 & [-0.577, -0.498] & 0.16 \\
\texttt{size\_na} & -0.757 & 0.030 & [-0.815, -0.708] & 0.111 \\
\texttt{txt\_na} & -0.439 & 0.087 & [-0.597, -0.298] & 0.28 \\
\texttt{bench\_na} & -0.094 & 0.054 & [-0.203, 0.003] & 0.76 \\
\texttt{lco2} & 0.842 & 0.097 & [0.601, 0.985] & 1.26 \\
\texttt{ds\_couse} & 1.121 & 0.019 & [1.081, 1.154] & 1.55 \\
\texttt{lquant} & 0.079 & 0.008 & [0.062, 0.091] & 1.23 \\
\texttt{lib\_transformers} & 0.203 & 0.033 & [0.146, 0.270] & 1.8 \\
\texttt{lib\_peft} & -1.801 & 0.169 & [-2.313, -1.559] & 0.00535 \\
\texttt{lib\_diffusers} & 0.594 & 0.153 & [0.246, 0.839] & 5.61 \\
\texttt{lib\_sentence-transformers} & 0.974 & 0.045 & [0.876, 1.072] & 16.9 \\
\texttt{redistributor} & 0.741 & 0.040 & [0.644, 0.813] & 8.61 \\
\midrule
\multicolumn{5}{l}{\textit{Quality}} \\
\texttt{bench\_z} & -0.116 & 0.025 & [-0.160, -0.059] & 0.971 \\
\midrule
\multicolumn{5}{l}{\textit{Demand match}} \\
\texttt{size\_dist} & -0.221 & 0.015 & [-0.246, -0.197] & 0.674 \\
\texttt{task\_hist} & 1.896 & 0.046 & [1.808, 1.973] & 246 \\
\texttt{lang\_hist} & 0.251 & 0.036 & [0.195, 0.308] & 2.07 \\
\texttt{ds\_lang} & 0.355 & 0.055 & [0.260, 0.459] & 2.8 \\
\texttt{ds\_task} & 0.883 & 0.109 & [0.746, 1.148] & 13 \\
\texttt{txt\_sim} & 1.687 & 2.078 & [-1.088, 4.529] & 1.15 \\
\midrule
\multicolumn{5}{l}{\textit{Inertia}} \\
\texttt{author\_family} & 0.052 & 0.059 & [-0.086, 0.146] & 1.16 \\
\texttt{same\_author} & 7.061 & 0.373 & [6.337, 7.891] & 8.06e+08 \\
\bottomrule
\end{tabular}
\end{table}


\paragraph{Reading the coefficients.}
Coefficients are utility weights. The last column of \cref{tab:coef} converts each into the factor on $(1+d)$ that
gives the same change in utility, using $\partial U / \partial \log d$ as the exchange rate; binary features are
switched from 0 to 1, continuous ones moved by one standard deviation. The two largest multiples belong to
\texttt{has\_deg} (a model with at least one earlier child) and \texttt{same\_author} (the author's own model): having
been adopted at all, and being one's own, dominate everything else.

\paragraph{Demand signals.}
Demand match can be detected only where the choice carries a demand signal: an author history, a declared dataset
with a usable profile, or at least five tokens of cleaned card text. \sigShare\% of first adoptions do. Restricted to
them, demand match accounts for \sigMatch\% of the gain; reweighting them to the composition of all first adoptions
(author experience $\times$ major-organization author) gives \sigMatchRw\%. The main estimate is a lower bound in
this sense: fit to needs that leave no trace in the data is attributed to no group.
\todo{Add: card-text threshold sensitivity (1/5/10/20/50 tokens); old-card robustness for 39 hubs.}

\section{Repeated Choices}
\label{app:repeat}

\paragraph{Decisions.}
Merging an author's consecutive uploads on one parent turns \nChoices{} choices into 226,542 / \nDecisions{} /
176,526 decisions for windows of 6 / 24 / 72 hours; \nFirst{} of them are first adoptions of an author--parent pair
and the rest are repeat decisions.

\paragraph{Stay or switch.}
Of the 75,016 first adoptions followed by another decision of the same author, 18.3\% stay with the parent at the
next decision, 5.8\% move to a newer version of the same family and 76.0\% switch to another model. A logit of
staying on the four group indices of the first parent (utility contributions under the main model), with controls
for the time to the next decision, whether the next task is the same, a newer family release and author experience,
gives positive dependence on the popularity index (0.056, SE 0.004) and negative dependence on the match index
($-0.165$, SE 0.017); a same-task next decision (0.95) and a newer family release ($-0.71$) matter most. Results are
similar for the same-task subsample, for staying within 180 days and with 6- and 72-hour windows.
\todo{Interpretation of the negative match coefficient; author-FE conditional logit; table.}

\paragraph{A two-nest model and why we do not rely on it.}
A natural model of repeated choice puts the author's earlier parents in an ``old'' nest and the rest of the market in
a ``new'' nest, with dissimilarity parameters $\lambda_{\text{old}}, \lambda_{\text{new}} \in (0,1]$. With a sampled new
nest the inclusive value must be estimated by Horvitz--Thompson, $\log \sum_{l \in \D} e^{V_l / \lambda} / \pi_l$
\citep{horvitz1952generalization}. On all \nDecisions{} decisions the estimates are $\hat\lambda_{\text{new}} =
0.021$ and $\hat\lambda_{\text{old}} = 0.128$, which would imply that the nest choice depends only on the best
alternative in each nest. A synthetic check with the same sampling design (\cref{tab:nestedcheck}) shows that
$\lambda$ is not reliably identified at the scale of the Hub. With a market of 50,000 the true $\lambda$ is recovered;
with 500,000 alternatives, close to the real median, a true $\lambda_{\text{new}} = 0.7$ is estimated at 0.10--0.47
with standard errors around 0.25, and enlarging the sampled new nest tenfold does not help. In a very large market
the new nest is chosen almost always, so the nest-level choice carries little information about $\lambda$. We
therefore report the binary stay-or-switch model and do not interpret $\hat\lambda$.

\begin{table}[h]
\centering
\caption{Synthetic recovery of $\lambda$ with a PPS-sampled new nest ($E|\D| = 300$ unless noted; 3,000 choices for
markets of 300 and 50,000, 1,500 for 500,000).}
\label{tab:nestedcheck}
\small
\begin{tabular}{rrrrr}
\toprule
Market & True $\lambda_{\text{new}}$ & Sample & $\hat\lambda_{\text{new}}$ (SE) & $\hat\lambda_{\text{old}}$ (SE) \\
\midrule
50,000 & 0.7 & 300 & 0.77 (0.12), 0.70 (0.15) & 0.64 (0.19), 0.22 (0.15) \\
50,000 & 0.3 & 300 & 0.36 (0.02), 0.33 (0.02) & 0.40 (0.05), 0.36 (0.05) \\
500,000 & 0.7 & 300 & 0.10 (0.25), 0.47 (0.26) & 0.15 (0.43), 0.11 (0.20) \\
500,000 & 0.7 & 3,000 & 0.42 (0.26) & 0.09 (0.18) \\
500,000 & 0.3 & 300 & 0.36 (0.04) & 0.36 (0.08) \\
\bottomrule
\end{tabular}
\end{table}
\todo{Generate this table from results/tables/v2/v2\_nested\_check.csv.}


\end{document}
